\begin{proposition}[Conductor topology over a higher-genus ruling]\hypertarget{tag-02YF}{}\label{night:P4R.3}\label{stmt:P4R.3}
\StatementRecord{night:P4R.3}{proved}{02YF}{\SeriesRef{cdp:st:A.5.3}, \SeriesRef{night:P4.12}, \SeriesRef{night:P4.13}, \SeriesRef{night:P4R.1}}

In the standing torus-fibration setting of this section, let an irreducible singular fibre
support $S$ have smooth normalization $\nu:T\to S$ whose
minimal model is ruled over a curve of genus $g\ge2$.
Write $D$ for the conductor divisor on $T$ and
$C=\operatorname{Sing}(S)_{\rm red}$.
The inclusion $D_{\rm red}\to T$ induces an isomorphism on
fundamental groups, and the normalization pushout induces
$\pi_1(S)\simeq\pi_1(C)$. Moreover,
\[
 e(C)\ge0,\qquad e(S)>0.
\]
The assertion includes conductor components on different
ruling fibres that map onto the same curve in $C$.
\end{proposition}
\begin{proof}

\begin{enumerate}[label={$\langle1\rangle$\arabic*.},leftmargin=*]
\item Proposition~\ref{cdp:st:A.5.3}, with the identity as the
minimal resolution of the smooth surface $T$, gives
$\omega_T(D)\simeq\nu^*\omega_S$ with $D$ effective and with its
conductor multiplicities. Proposition~\ref{night:P4.12} makes this
line torsion, so $K_T+D\equiv0$.
Let $\pi:T\to T_0$ be the sequence of point blowdowns to the
minimal ruled surface $r_0:T_0\to B$, and put $D_0=\pi_*D$.
By \cite[III.10, p.~40]{Beauville1978}, write $T_0=\mathbf P_B(V)$
for a rank-two bundle $V$ on $B$. In the basis of the tautological
class $h=c_1(\OO_{\mathbf P_B(V)}(1))$ and the fibre class $F$,
$h^2=\deg V$, $h\cdot F=1$, $F^2=0$ and
$K_{T_0}=-2h+(\deg V+2g-2)F$
\cite[III.18, pp.~46--47]{Beauville1978}.
Thus $K_{T_0}+D_0\equiv0$ and $D_0\cdot F=2$.

\item There are three possibilities for the horizontal part of
$D_0$: a reduced bisection, two distinct reduced sections,
or twice one section. The first two are excluded directly.
For an irreducible horizontal component $H$ of coefficient
one, adjunction and the effectiveness of the other
components give
\[
 2p_a(H)-2=(K_{T_0}+H)\cdot H
       =-(D_0-H)\cdot H\le0.
\]
On the normalization of $H$, Riemann--Hurwitz for its
finite map of degree $d\ge1$ to $B$ gives
$2g(\widetilde H)-2\ge d(2g-2)>0$.
Since $p_a(H)\ge g(\widetilde H)$, these inequalities
are incompatible. Thus
\[
 D_0=2H+\sum_{j=1}^{v}a_jF_j,\qquad a_j\ge1,
\]
where $H$ is a section and the $F_j$ are distinct ruling
fibres. A horizontal curve of degree one is a section:
its normalization maps isomorphically to the smooth base,
and the resulting map to the projective bundle is a
section with closed image. In particular $H$ is smooth.

\item At one step $\beta:T'\to T''$ of the blowup sequence,
let $\mu$ be the multiplicity of the conductor $D''$ at the
centre and $E$ the exceptional curve.
Write $D'=D''_{\rm strict}+aE$. The canonical blowup formula
\cite[II.3(iv), pp.~17--18]{Beauville1978} gives
\[
 K_{T'}+D'
   =\beta^*(K_{T''}+D'')+(1+a-\mu)E.
\]
Intersection with $E$ and numerical triviality give
$a=\mu-1$. Effectiveness therefore places every centre
on the conductor. At a smooth conductor point of
coefficient one, the exceptional curve is absent from
the new conductor. At a smooth point of coefficient at
least two, it is a rational leaf attached at one point.
At an intersection of two conductor components it
subdivides their edge. Starting with $H$ and the
transverse fibres $F_j$, these operations preserve
normal crossings and make $D_{\rm red}$ a copy of $H$
with rational trees attached. No cycle is introduced.

\item The projection of $H$ to $B$ induces an isomorphism on
fundamental groups, and every attached rational tree
is simply connected. A point blowup replaces a four-ball by the disc bundle
of $\OO_{\mathbb P^1}(-1)$, with common boundary $S^3$.
Van Kampen therefore shows that it leaves $\pi_1(T)$ unchanged. The preceding
description also proves, at every blowup, that the
inclusion of the reduced conductor induces the same
isomorphism to the fundamental group of the surface.
Hence $\pi_1(D_{\rm red})\to\pi_1(T)$ is an
isomorphism.

\item The normalization morphism $\nu$ is finite, is an
isomorphism away from the conductor, and has smooth
source. Its nonisomorphism locus equals $C$: a normal
point has a neighbourhood identified with its
normalization and is consequently smooth. Conversely
every smooth point is normal. Thus the natural map
\[
 T\underset{D_{\rm red}}{\cup}C\longrightarrow S
\]
is a continuous bijection. Both source and target are
compact Hausdorff spaces, so it is a homeomorphism.
A triangulation compatible with the analytic pair $(T,D_{\rm red})$
and the finite curve map $D_{\rm red}\to C$ makes
$D_{\rm red}\subset T$ a subcomplex and hence a
cofibration. Replacing the map $D_{\rm red}\to C$ by
its mapping cylinder computes the same homotopy
pushout: the homotopy extension property for this
cofibration contracts the inserted cylinder without
altering the homotopy type. Van Kampen on that double
mapping cylinder therefore gives
\[
 \pi_1(S)=\pi_1(T)
       *_{\pi_1(D_{\rm red})}\pi_1(C)
       \simeq\pi_1(C).
\]
All three spaces in the amalgam are connected. This
argument includes every identification made by the
finite conductor map, including identifications
between horizontal and vertical branches.

\item Proposition~\ref{night:P4R.1} makes this group virtually abelian. Let $c$ be
the number of irreducible components of $C$, let
$g_i$ be the genera of their smooth normalizations,
and let $b$ be the first Betti number of the incidence
graph of the finite point identifications. Replacing
each identification by a small tree joining its
preimages gives a homotopy model consisting of these
closed oriented surfaces joined along a connected
graph. Choosing a spanning tree gives
\[
 \pi_1(C)\simeq
   \Bigl(*_{i=1}^{c}\pi_1(C_i^{\rm norm})\Bigr)*F_b,
 \qquad
 e(C)=c+1-2\sum_i g_i-b.
\]
The formula also covers unibranch singular points:
their normalization is a local topological
homeomorphism and they contribute no graph edge.

\item A genus-$h$ surface factor with $h\ge2$ cannot be virtually
abelian: a connected degree-$d$ cover has genus $1+d(h-1)>1$,
and its fundamental group is nonabelian. Projection onto this
surface factor is a quotient map, while every quotient of a
virtually abelian group is virtually abelian. Two nontrivial surface factors,
two free generators, or one elliptic factor together
with a free generator contain a nonabelian free
subgroup. Such a subgroup cannot lie in a virtually
abelian group: its intersection with a finite-index
abelian subgroup would have finite index and would
be abelian, whereas a finite-index subgroup of a
nonabelian free group has rank at least two. The
possibilities are therefore all $g_i=0$ with
$b\le1$, or exactly one $g_i=1$ with $b=0$.
The displayed Euler formula gives $e(C)\ge0$.
Additivity of compactly supported Euler numbers and
Proposition~\ref{night:P4.13} give
\[
 e(S)=e(S_{\rm reg})+e(C)\ge0.
\]

\item Finally equality is impossible here. It would force
$e(C)=0$, so the same formula gives $c=1$, $g_1=1$
and $b=0$. Every irreducible component of
$D_{\rm red}$ maps finitely and nonconstantly onto
this curve. A rational vertical component would
therefore give a nonconstant map from $\mathbb P^1$
to its genus-one normalization, contrary to
Riemann--Hurwitz. There are no vertical conductor
components. The initial conductor is consequently
$2H$, and there can have been no blowup: the first
centre would lie on $2H$ and would create a vertical
conductor component of coefficient one. Thus
$T=T_0$ and $D_{\rm red}=H$. It follows that
\[
 e(S_{\rm reg})=e(T\setminus H)
       =(4-4g)-(2-2g)=2-2g<0,
\]
contradicting Proposition~\ref{night:P4.13}. This proves strict positivity.

\end{enumerate}
\end{proof}